\chapter{Complexity — Meyer Machine Walk}
%%% generated by the blueprint pipeline from TCSlib.Complexity.CircuitComplexity.MeyerMachineWalk
\section{Overview}

Declarations of the Lean module \texttt{TCSlib.Complexity.CircuitComplexity.MeyerMachineWalk}.

\section{Declarations}

\begin{definition}[Configuration Update from a Selector]
\lean{Complexity.Meyer.wstep}
\label{Complexity.Meyer.wstep}
\leanok
\uses{Complexity.Meyer.Sel, Complexity.Meyer.walkDir, Turing.Cfg, Turing.moveInputPos}
Let $M$ be a finite Turing machine over the Boolean alphabet, let $x$ be a Boolean
string, and let $c$ be a configuration of $M$ on input $x$. A selector specifying a work
tape $h$ and direction $d$ changes only that work tape's head position, by adding $d$; a
selector specifying the input tape and direction $d$ changes only the input-head
position by one step in direction $d$. A selector specifying a state leaves the
configuration unchanged.
\end{definition}

\begin{theorem}[Output when the decision tape is blank]
\lean{Complexity.Meyer.wdec_emit}
\label{Complexity.Meyer.wdec_emit}
\leanok
\uses{Complexity.Meyer.aux_read0, Complexity.Meyer.aux_read2, Complexity.Meyer.decodeSel, Complexity.Meyer.lcfg, Complexity.Meyer.lcfg_apply_aux, Complexity.Meyer.lcfg_step, Complexity.Meyer.mreads, Complexity.Meyer.selAnswer, Complexity.Meyer.tabTM, Complexity.Meyer.tabTr, Complexity.Meyer.tcfg, Complexity.TimeHierarchy.OutReg, Turing.Cfg.inputSymbol, Turing.Cfg.workTapeSymbols, Turing.MultiTapeTM.step}
\difficulty{3}
Let $M$ be a finite-state Turing machine with Boolean symbols, let $x,y$ be Boolean
words, let $bits$ be a Boolean selector encoding, let $P$ be a valid input-head position
for $y$, and let $c$ be a configuration of $M$ on $x$. Given auxiliary tape contents
$\tau_C,\tau_O$, positions $z_C,z_O$, an optional state $oq$, an output register $r$,
and a Boolean tag $b$, if $\tau_O(z_O)$ is blank, then one step of the constructed
machine from its waiting-decision configuration halts and outputs the single bit
answering the decoded selector from $oq$, $r$, and the symbols currently scanned by $c$.
\end{theorem}

\begin{theorem}[Zero-bit decision step]
\lean{Complexity.Meyer.wdec_zero}
\label{Complexity.Meyer.wdec_zero}
\leanok
\uses{Complexity.Meyer.aux_read2, Complexity.Meyer.lcfg, Complexity.Meyer.lcfg_apply_aux, Complexity.Meyer.lcfg_step, Complexity.Meyer.tabTM, Complexity.Meyer.tabTr, Complexity.TimeHierarchy.OutReg, Turing.MultiTapeTM.step}
\difficulty{3}
Let $M$ be a finite-state Turing machine with Boolean symbols, let $x,y$ be Boolean
words, let $bits$ be a Boolean selector encoding, let $P$ be a valid input-head position
for $y$, and let $c$ be a configuration of $M$ on $x$. Given auxiliary tape contents
$\tau_C,\tau_O$, positions $z_C,z_O$, an optional state $oq$, an output register $r$,
and a Boolean tag $b$, if $\tau_O(z_O)$ contains false, then one step of the constructed
machine leaves it in the waiting-decision state, changes $\tau_O(z_O)$ to true, and
moves that auxiliary head to $z_O+1$, with the other listed configuration data
unchanged.
\end{theorem}

\begin{theorem}[Output-tape return step at a symbol]
\lean{Complexity.Meyer.wret_some}
\label{Complexity.Meyer.wret_some}
\leanok
\uses{Complexity.Meyer.aux_read2, Complexity.Meyer.lcfg, Complexity.Meyer.lcfg_apply_aux, Complexity.Meyer.lcfg_step, Complexity.Meyer.tabTM, Complexity.Meyer.tabTr, Complexity.TimeHierarchy.OutReg, Turing.MultiTapeTM.step}
\difficulty{3}
Let $M$ be a finite-state Boolean Turing machine, let $x,y$ be Boolean words, and fix an
encoded selector, a simulated configuration $c$ on input $x$, auxiliary tape contents
and head positions, a position $P$ on input $y$, an output register $r$, and a Boolean
tag $b$. If the output tape contains the bit $v$ at its current head position $z_O$,
then one step of the tabulating machine from the lifted state with control state
$\mathrm{wret}$ leaves the other data and output tape unchanged, keeps the control state
at $\mathrm{wret}$, and moves the output-tape head from $z_O$ to $z_O-1$.
\end{theorem}

\begin{theorem}[Output-tape return step at a blank cell]
\lean{Complexity.Meyer.wret_none}
\label{Complexity.Meyer.wret_none}
\leanok
\uses{Complexity.Meyer.aux_read2, Complexity.Meyer.lcfg, Complexity.Meyer.lcfg_apply_aux, Complexity.Meyer.lcfg_step, Complexity.Meyer.tabTM, Complexity.Meyer.tabTr, Complexity.TimeHierarchy.OutReg, Turing.MultiTapeTM.step}
\difficulty{3}
Let $M$ be a finite-state Boolean Turing machine, let $x,y$ be Boolean words, and fix an
encoded selector, a simulated configuration $c$ on input $x$, auxiliary tape contents
and head positions, a position $P$ on input $y$, an output register $r$, and a Boolean
tag $b$. If the output tape is blank at its current head position $z_O$, then one step
of the tabulating machine from the lifted state with control state $\mathrm{wret}$
leaves the other data and output tape unchanged, changes the control state to
$\mathrm{wdec}$, and moves the output-tape head from $z_O$ to $z_O+1$.
\end{theorem}

\begin{theorem}[Firing the decision sweep]
\lean{Complexity.Meyer.wdec_fire}
\label{Complexity.Meyer.wdec_fire}
\leanok
\uses{Complexity.Meyer.aux_read0, Complexity.Meyer.aux_read2, Complexity.Meyer.decodeSel, Complexity.Meyer.idle, Complexity.Meyer.lcfg, Complexity.Meyer.lcfg_apply_aux, Complexity.Meyer.lcfg_step, Complexity.Meyer.one, Complexity.Meyer.tabTM, Complexity.Meyer.tabTr, Complexity.Meyer.tact_apply, Complexity.Meyer.walkDir, Complexity.Meyer.wstep, Complexity.TimeHierarchy.OutReg, Turing.Cfg.inputSymbol, Turing.FinTM.VirtualTag, Turing.FinTM.virtualMove, Turing.FinTM.virtualMove_correct, Turing.FinTM.virtualNextTag, Turing.MultiTapeTM.step, Turing.moveInputPos_zero}
\difficulty{6}
Let $M$ be a finite-state Turing machine with Boolean symbols, let $x,y$ be Boolean
words, let $bits$ be a Boolean selector encoding, let $P$ be a valid input-head position
for $y$, and let $c$ be a configuration of $M$ on $x$. Given auxiliary tape contents
$\tau_C,\tau_O$, positions $z_C,z_O$, an optional state $oq$, an output register $r$,
and a Boolean tag $b$ that correctly marks the boundary position of $c$'s input head, if
$\tau_O(z_O)$ contains true, then there is a tag $b'$ correctly marking the boundary
position after one decoded-selector walk step on $c$ such that one step of the
constructed machine gives the waiting-return configuration with that walked
configuration, with $\tau_O(z_O)$ changed to false and the auxiliary head moved to
$z_O-1$.
\end{theorem}

\begin{theorem}[Sweep across a Boolean word]
\lean{Complexity.Meyer.wdec_sweep}
\label{Complexity.Meyer.wdec_sweep}
\leanok
\uses{Complexity.Meyer.bufferTape_at, Complexity.Meyer.lcfg, Complexity.Meyer.tabTM, Complexity.Meyer.wdec_zero, Complexity.TimeHierarchy.OutReg, Complexity.TimeHierarchy.bufferTape_update_mid, Turing.Cfg, Turing.FinTM.bufferTape, Turing.MultiTapeTM.runFrom, Turing.MultiTapeTM.runFrom_succ_eq_step}
\difficulty{5}
Let $M$ be a finite-state Turing machine with Boolean symbols, let $x,y$ be Boolean
words, let $bits$ be a Boolean selector encoding, let $P$ be a valid input-head position
for $y$, and let $c$ be a configuration of $M$ on $x$. For any optional state $oq$,
output register $r$, Boolean tag $b$, auxiliary tape content $\tau_C$, integer position
$z_C$, Boolean word $rest$, and natural numbers $i,j$, running the constructed machine
for $j$ steps from its waiting-decision configuration with output tape content
consisting of $i$ true bits, then $j$ false bits, then $rest$, and auxiliary head at
$i$, yields the same waiting-decision configuration with the first $i+j$ bits true,
followed by $rest$, and auxiliary head at $i+j$.
\end{theorem}

\begin{theorem}[Output-tape return sweep]
\lean{Complexity.Meyer.wret_sweep}
\label{Complexity.Meyer.wret_sweep}
\leanok
\uses{Complexity.Meyer.bufferTape_at, Complexity.Meyer.lcfg, Complexity.Meyer.tabTM, Complexity.Meyer.wret_none, Complexity.Meyer.wret_some, Complexity.TimeHierarchy.OutReg, Turing.Cfg, Turing.FinTM.bufferTape, Turing.MultiTapeTM.runFrom, Turing.MultiTapeTM.runFrom_succ_eq_step, Turing.MultiTapeTM.runFrom_zero}
\difficulty{4}
Let $M$ be a finite-state Boolean Turing machine, let $x,y$ be Boolean words, and fix
encoded selector bits, a simulated configuration $c$ on input $x$, auxiliary tape
contents and head positions, an input position $P$ on $y$, an optional simulated state
$oq$, an output register $r$, a Boolean tag $b$, and a Boolean word $L$. For every
natural number $j$, starting in control state $\mathrm{wret}$ with the output tape
storing $j$ ones followed by $L$ and its head at position $j-1$, after $j+1$ steps the
control state is $\mathrm{wdec}$, the tape contents are unchanged, and its head is at
position $0$.
\end{theorem}

\begin{theorem}[Bounded simulation of a walk]
\lean{Complexity.Meyer.walk_run}
\label{Complexity.Meyer.walk_run}
\leanok
\uses{Complexity.Meyer.bufferTape_at, Complexity.Meyer.decodeSel, Complexity.Meyer.lcfg, Complexity.Meyer.selAnswer, Complexity.Meyer.tabTM, Complexity.Meyer.wdec_emit, Complexity.Meyer.wdec_fire, Complexity.Meyer.wdec_sweep, Complexity.Meyer.wret_sweep, Complexity.Meyer.wstep, Complexity.TimeHierarchy.OutReg, Complexity.TimeHierarchy.bufferTape_update_mid, Complexity.TimeHierarchy.ctrVal, Complexity.TimeHierarchy.ctrVal_borrow, Complexity.TimeHierarchy.ctrVal_pos_decomp, Complexity.TimeHierarchy.ctrVal_zero_eq, Turing.Cfg, Turing.Cfg.inputSymbol, Turing.Cfg.workTapeSymbols, Turing.FinTM.VirtualTag, Turing.FinTM.bufferTape, Turing.MultiTapeTM.runFrom, Turing.MultiTapeTM.runFrom_add, Turing.MultiTapeTM.runFrom_succ_eq_step'}
\difficulty{6}
Let $M$ be a finite-state Turing machine with Boolean input and work-tape symbols, let
$x,y$ be Boolean words, let $bits$ be a Boolean encoding of the selectors of $M$, and
let $P$ be an input-head position for $y$. For any proposed output state $oq$, output
register $r$, auxiliary tape content $\tau_C$ and position $z_C$, natural number $v$,
Boolean word $w$, configuration $c$ of $M$ on $x$, and tag $b$, if $w$ has counter value
$v$ and $b$ correctly marks the boundary position of $c$'s input head, then the machine
constructed from $M$, started in its waiting-decision state with these data, halts
within $(v+1)(2|w|+3)$ steps and outputs the single bit that answers the decoded
selector using $oq$, $r$, and the symbols read from $c$ after $v$ simulated
selector-walk steps.
\end{theorem}
